\documentclass{article}
\usepackage[a4paper]{geometry}
\usepackage{fancyhdr}
\pagestyle{fancy}
\lhead{Globalisierung}
\rhead{Februar 2026}
\begin{document}
Die Globalisierung beschreibt die Weltweite wachsende wirtschaftliche Vernetzung, den kulturellen Austausch und Regierungsbündnisse.
 
Hier wird sich auf den wirtschaftlichen Aspekt fokusiert. 
 
\section{Triebkräfte und Indikatoren}
Die wirtschaftliche Globalisierung wurde durch unterschiedlichste Triebkräfte vorangetrieben.
\begin{description}
 \item[Deregulierung] Der Abbau von der Bürokratie eines Staates, welche den globalen Handel bisher teurer und langsamer machte.
 \item[Liberalisierung] Bilden eines freieren globalen Marktes z.\,B. durch das Abbauen von Zöllen, bilden von Freihandelsabkommen.
 \item[Privatisierung] Bisherig staatliche Organisationen dem Freien Mark übergeben, in der Hoffnung, dass dieser Effizienter sei.
 \item[Technologie] Technologischer Fortschritt erlaubte überhaupt den globalen Handel, z.\,B. durch die Erfindung des effizienten Transportes über Seefahrt und Flugzeuge, oder der Erfindung der effizienten Kommunikation über Telefone oder das Internet.
\end{description}  
Die wirtschaftliche Globalisierung zeigt unterschiedlichste Indikatoren auf. Dazu zählen ein steigendes (globales) BIP, steigende globale Export und Importe, steigende Investitionen in das Ausland und steigende Anteile einer Währung am Devisenmarkt.
\end{document}